\section{Decimonovena iteración de la Fase 5: Carga unificada de cartuchos}
\label{sec:F5I19}

En esta iteración abordamos unificar las cargas de ficheros para el manejo de los cartuchos: una nueva pestaña para la carga de cartucho sustituye a los botones de carga repartidos por el resto de pestañas. Al cargar un cartucho, como se ha incluido la información de los anteriores en el sistema ya no es necesario cargar en el orden predefinido: GWS $\to$ GWB $\to$ GWI, cualquier cartucho nos dará los datos de la pestaña Setup. No cambia el núcleo de carga (F5-I9): cambia la puerta.

La excepción es que para los cartuchos GWC es necesario cargar antes uno que incluya los datos del setup, pues los datos básicos de pantalla y botones no viajan en el cartucho GWC.

\subsection{Análisis del estado actual}
\label{sec:F5I19Analisis}

\begin{itemize}
	\item El núcleo ya está unificado (F5-I9): \texttt{\_cargarCartucho(ficheros,
		tipoEsperado)} decide la forma por el CONTENIDO (ZIP $\to$ descomprime,
	\texttt{.gwc} $\to$ flujo binario, resto $\to$ flujo de carpeta) y hay un
	mensaje de error único.
	\item La puerta no: cinco \texttt{input} ocultos y ocho botones de carga
	en tres vistas, encuadrados por \texttt{BOTONES\_ACCION}; los botones de
	cartucho exigen setup aplicado y la lectura valida contra el estilo del
	setup (\texttt{setupGWS.estilo}).
	\item Sobran, y se retiran aquí (R07-B del I99): exports muertos
	\texttt{leerGWB}, \texttt{leerGWI}, \texttt{leerJSON}, \texttt{escalar},
	\texttt{indiceDe}, \texttt{bloque}; \texttt{cargarDelMapa};
	\texttt{'GWC'} en \texttt{TIPOS\_CARTUCHO}; \texttt{json.img\_s} en la
	raíz.
\end{itemize}

\subsection{Diseño: la pestaña de carga}
\label{sec:F5I19Diseno}

\begin{itemize}
	\item \textbf{Una pestaña nueva, siempre activa.} En la barra de
	pestañas, visible y habilitada como las demás, salvo durante la partida
	(la función \texttt{jugar()} ya oculta las pestañas). Sin botón que la
	abra: es una pestaña más.
	\item \textbf{El panel admite todo.} Dentro:
	\begin{enumerate}
		\item botón \textbf{Cargar ficheros} — \texttt{input} \texttt{multiple} para \texttt{.json} (GWS), \texttt{.zip} (GWB/GWI) y \texttt{.gwc}; 
		\item botón \textbf{Cargar carpeta} — el \texttt{input} \texttt{webkitdirectory}; 
		\item zona de arrastre que acepta ficheros o carpetas y el pegado con Ctrl+V donde el navegador dé \texttt{clipboardData} (con archivos; lo demás, aviso).
	\end{enumerate}	 
	\item \textbf{El tipo lo decide el contenido} (detector de forma F5-I9):
	adiós al \texttt{tipoEsperado} de los botones. Un GWS pobla la pestaña
	Setup; un GWB o un GWI poblan \textbf{ambas} pestañas (autoría y setup)
	y construyen la máquina; un GWC solo juega con la máquina ya construida.
	\item \textbf{El cartucho cargado manda.} En cada carga de GWB/GWI la
	máquina se re-ajusta al estilo del cartucho y las pestañas Setup y
	Cartucho se poblan con sus datos (teclas incluidas; los \texttt{puerto}
	del ESP32 quedan vacíos si no hay setup previo, pues no viajan en el
	cartucho). \texttt{\#estado} avisa cuando el setup mostrado es derivado
	del cartucho y no un GWS cargado.
	\item \textbf{GWC con condición.} Como el binario no viaja ni estilo ni
	botonera, exige máquina previa: de un setup o derivada de un GWB/GWI ya
	cargado; sin ella, error claro: «falta la configuración de la
	máquina, defínala o cargue un cartucho que la incluya (F5-I19)».
	\item \textbf{GWB autosuficiente al generarse}. El GWB que guarda el
	programa declara siempre \texttt{estilo}, \texttt{botones} (con sus
	\texttt{tecla}) y resolución, de modo que un GWB guardado carga sin
	setup. Un GWB viejo que declare menos requiere setup previo y da un error
	claro que lo diga.
	\item \textbf{Adiós a las copias.} Desaparecen los ocho botones de carga
	y los inputs \texttt{ficheroSetup}, \texttt{ficheroGWBZip},
	\texttt{ficheroGWIZip} y \texttt{ficheroGWC}, sustituidos por los dos
	\texttt{input} del panel (\texttt{ficheroCartucho} y
	\texttt{carpetaCartucho}); las
	listas \texttt{BOTONES\_ACCION}/\texttt{BOTONES\_CARTUCHO} se recortan a
	lo que quede vivo. Las acciones Aplicar y Guardar conservan sus botones.
\end{itemize}

\subsection{Plan de cambios}
\label{sec:F5I19Plan}

\begin{itemize}
	\item \texttt{index.html}: pestaña y panel nuevos; fuera los botones e
	inputs duplicados.
	\item \texttt{BaseDirector.js}: escuchadores del panel (change de ambos
	inputs, dragover/drop, paste) todos hacia
	\texttt{\_cargarCartucho(ficheros)} sin tipo esperado;
	\texttt{validarComunes} toma la condición del GWB sin estilo (error
	claro); la reconstrucción de máquina pobla ambas pestañas; recorte de
	\texttt{BOTONES\_ACCION}/\texttt{BOTONES\_CARTUCHO}; cabeceras al día.
	\item \texttt{BaseCartucho.js}: retirar restos R07-B; exclusión de
	\texttt{'GWC'} de \texttt{TIPOS\_CARTUCHO}; cabeceras al día.
\end{itemize}

\subsection{Pruebas}
\label{sec:F5I19Pruebas}

\begin{itemize}
	\item Las suites de pruebas S2, S3 y S5 se deben ejecutar sin problemas y deben seguir en verde; sus cajas usan
	\texttt{U.cargarEn('ficheroCartucho',\ldots)} (los ids antiguos
	\texttt{ficheroGWBZip}, \texttt{ficheroGWIZip} y \texttt{ficheroGWC}
	desaparecen), así que los
	arrancadores se re-apuntan a los inputs nuevos del panel (cambio
	mecánico, misma prueba). Las referencias de F5-I18 se actualizan si
	cambia el id.
	\item Cajas nuevas (grupo \texttt{director}; D5 no es destructiva y
	DEBE SER LA PRIMERA de la suite S2, pues necesita la página recién
	cargada; D6 y D7 sustituyen el cartucho) que se
	registran y se REFERENCIAN EN F5-I18, la norma vigente:
	(1) \textbf{carga sin setup} (caja D6, tras la D5): GWB con
	\texttt{estilo} declarado y GWI
	cargan y poblan Setup + Cartucho; \texttt{\#estado} avisa «derivado»;
	(2) \textbf{GWC sin máquina} (caja D5, la PRIMERA de la suite S2):
	error claro; (3) \textbf{cambio de
		máquina al recargar} (caja D7): tras un setup2, cargar un GWB de
	estilo 4 y
	comprobar que se ajusta todo (estilo, botonera, lienzo, sprites y
	título) — es el punto pendiente de comprobar al re-ajustar la máquina
	en caliente.
	\item Al cerrar la iteración: S7 completa en verde.
\end{itemize}

\subsection{Criterios de aceptación}
\label{sec:F5I19Criterios}

\begin{enumerate}
	\item Una sola pestaña carga GWS, GWB (carpeta o ZIP), GWI y GWC, desde
	cualquier estado, sin orden predefinido.
	\item El cartucho cargado manda: Setup y Cartucho quedan poblados con
	sus datos y la máquina re-ajustada (con la caja~3 verde).
	\item El GWC exige máquina previa, con error claro.
	\item Un GWB generado por el programa carga solo, sin setup.
	\item S1--S6 en verde; S7 completa al cerrar.
\end{enumerate}